\documentclass[a4paper]{article}

\usepackage[spanish]{babel}
\usepackage{graphicx}
\usepackage{amsmath, amssymb}
\usepackage[margin=2cm]{geometry}
\usepackage{fancyhdr}
\usepackage{enumerate}
\usepackage[shortlabels]{enumitem}
\usepackage{parskip}
\usepackage[most]{tcolorbox}
\usepackage[hidelinks]{hyperref}
\usepackage{float}
\usepackage{booktabs}
\usepackage{mathrsfs}
\usepackage{tikz}

% cabecera
\pagestyle{fancy}
\fancyhead[l]{Astroestadística}
\fancyhead[c]{Quiz 3 ???}
\fancyhead[r]{\today}
\fancyfoot[c]{\thepage}
\renewcommand{\headrulewidth}{0.2pt} % linea horizontal


\begin{document}

\section{Quiz de Pruebas de Hipótesis}

Elija la respuesta correcta en cada caso.

\begin{enumerate}

\item Una prueba bilateral se utiliza cuando:
\begin{enumerate}[a)]
    \item Se desconoce la desviación estándar poblacional.
    \item Se desea detectar cualquier desviación, sea por exceso o defecto.
    \item Se quiere saber si el parámetro es mayor que un valor específico.
    \item Se quiere confirmar si el parámetro es menor que un valor específico.
\end{enumerate}

\item ¿Por qué es importante definir la hipótesis alternativa antes de conocer los resultados de los datos?
\begin{enumerate}[a)]
    \item Porque así se puede suponer normalidad.
    \item Porque lo exige el software estadístico.
    \item Para facilitar los cálculos.
    \item Para evitar sesgos y manipulación de los resultados.
\end{enumerate}

\item ¿Cuál de las siguientes afirmaciones describe correctamente el error tipo II en una prueba de hipótesis?
\begin{enumerate}[a)]
    \item No rechazar $H_0$ cuando $H_0$ es falsa.
    \item Aceptar $H_0$ cuando $H_0$ es verdadera.
    \item Rechazar $H_a$ cuando $H_a$ es verdadera.
    \item Rechazar $H_0$ cuando $H_0$ es verdadera.
\end{enumerate}

\item ¿Qué efecto tiene reducir el nivel de significancia $\alpha$ (por ejemplo, de 0.05 a 0.01) en la probabilidad del error tipo II?
\begin{enumerate}[a)]
    \item Disminuye.
    \item No cambia.
    \item Aumenta.
    \item Elimina el error tipo II.
\end{enumerate}

\item En una prueba de hipótesis, el valor-p representa:
\begin{enumerate}[a)]
    \item La probabilidad de encontrar un resultado igual o más acorde a $H_0$ que el observado.
    \item La probabilidad de obtener el valor observado, suponiendo que $H_a$ es cierta.
    \item La probabilidad de obtener un resultado igual o más contradictorio a $H_0$ que el observado.
    \item La probabilidad de cometer un error tipo I.
\end{enumerate}

\end{enumerate}

\bibliographystyle{plainurl}
\bibliography{bibliografia} 
\end{document}
